\section{为什么要 baseline}

\subsection{高方差问题}

在语言模型 RL 中，reward 往往是稀疏的、延迟的、并且分布差异很大。字幕里的 toy 例子把这个问题讲得很直白：

\begin{itemize}
    \item state $s_1$ 的 reward 可能是 11 或 9；
    \item state $s_2$ 的 reward 可能是 0 或 2；
    \item 直接拿 reward 做更新，会出现“9 比 2 大，所以看起来更值得强化”的局部错觉。
\end{itemize}

但是从全局看，$s_2$ 明显更难、reward 更低，不能只看绝对值。这里真正要比较的是“\textbf{这个 action 比在当前 state 下的平均水平好多少}”。

\subsection{baseline 不改变目标}

引入 baseline $b(s)$ 后，我们优化的是
\[
\mathbb{E}[R - b(s)].
\]
只要 baseline 不依赖动作 $a$，这个变换不会改变原来的最优解，因为
\[
\mathbb{E}[b(s)] = \int p(s)\,b(s)\,ds
\]
与 policy 无关，因而只是加了一个常数项。

\begin{knowledgebox}{baseline 的要求}
baseline 可以依赖 state，不能依赖 action。直观上，它是“当前 state 下你本来就该得到多少”的参照线。
\end{knowledgebox}

\begin{table}[H]
\centering
\begin{tabular}{@{}cccc@{}}
\toprule
\textbf{state} & \textbf{action 1} & \textbf{action 2} & \textbf{直觉} \\
\midrule
$s_1$ & 11 & 9 & 两个都还不错，但 11 更好 \\
$s_2$ & 0 & 2 & 2 仍然优于 0，但绝对值更低 \\
\bottomrule
\end{tabular}
\caption{字幕里的两状态例子：关键不在绝对 reward，而在当前 state 下的相对提升}
\end{table}

\begin{knowledgebox}{baseline 的数学意义}
如果 baseline 只依赖 state，那么
\[
\mathbb{E}_{a \sim \pi(\cdot \mid s)}[\nabla \log \pi(a \mid s)\,b(s)] = 0.
\]
因此我们可以安全地把 reward 改成 $R-b(s)$，而不改变期望梯度的方向。
\end{knowledgebox}


\subsection{方差为什么会下降}

在 toy 例子里，原始 reward 的方差大约是 5.3；如果对两个 state 分别取 baseline 10 和 1，方差会降到约 1.1。这个例子想说明的不是某个数值，而是一个原则：

\begin{importantbox}{baseline 的作用}
\begin{itemize}
    \item 保留期望优化方向；
    \item 压缩无用的绝对量级差异；
    \item 让梯度更新更稳定、更快收敛。
\end{itemize}
\end{importantbox}


\subsection{advantage function}

经典 RL 里有：
\[
V(s) = \mathbb{E}[R \mid s], \qquad Q(s,a)=\mathbb{E}[R \mid s,a],
\]
以及 advantage：
\[
A(s,a)=Q(s,a)-V(s).
\]

在这一讲的 outcome reward 设定里，$Q$ 和 $R$ 基本可以视为同一个量，因此当 baseline 取成
\[
b(s) = \mathbb{E}[R \mid s]
\]
时，$R - b(s)$ 就是 advantage 的形式。

\begin{warningbox}{难点不在公式，而在估计}
理论上最优 baseline 可以写出 closed form，但实际里很难精确计算期待值，所以大家通常用样本平均、group mean，或者 value model 去近似它。
\end{warningbox}


